%!TEX program = xelatex
%!TEX encoding = UTF-8
\def\ano{2024}
\def\semestre{1}
\def\disciplina{Álgebra 1}
\def\turma{C}
\def\autor{José Antônio O. Freitas}
\def\instituto{MAT-UnB}

\documentclass{beamer}
\usetheme{Madrid}
\usecolortheme{beaver}
% \mode<presentation>
\usepackage[singlelinecheck=false]{caption}
\usepackage{lmodern}
\usepackage{textpos}
\usepackage{amssymb}
\usepackage{amsmath,amsfonts,amsthm,amstext}
\usepackage{cases}
\usepackage[justification=centering]{caption}
\usepackage[brazilian]{babel}
% \usepackage[latin1]{inputenc}
\usepackage{graphicx}
\graphicspath{{/home/jfreitas/GitHub/geometria_analitica/apresentacoes/logo}{/home/jfreitas/GitHub/geometria_analitica/notas_de_aula/figuras-latex}{C:/Users/jfreitas/GitHub/geometria_analitica/apresentacoes/logo}}
\usepackage{enumitem}
\usepackage{multicol}
\usepackage{answers}
\usepackage{qrcode}
\usepackage{tikz,ifthen}
\usetikzlibrary{lindenmayersystems}
\usetikzlibrary[shadings]

\newcounter{saveenumi}
\newcommand{\seti}{\setcounter{saveenumi}{\value{enumi}}}
\newcommand{\conti}{\setcounter{enumi}{\value{saveenumi}}}

\resetcounteronoverlays{saveenumi}


\renewcommand{\qedsymbol}{$\blacksquare$}

\setbeamertemplate{itemize items}[ball]

\input{../configuracoes_latex/ambientes}
\newcommand{\nsub}{\varsubsetneq}
\newcommand{\vaz}{\emptyset}
\newcommand{\dom}{{\rm dom\,}}
\newcommand{\im}{{\rm Im\,}}
\newcommand{\sub}{\subseteq}
\newcommand{\n}{\mathbb{N}}
\newcommand{\z}{\mathbb{Z}}
\newcommand{\rac}{\mathbb{Q}}
\newcommand{\real}{\mathbb{R}}
\newcommand{\complex}{\mathbb{C}}
\newcommand{\cp}[1]{\mathbb{#1}}
\newcommand{\ch}{\mbox{\textrm{car\,}}\nobreak}
\newcommand{\vesp}[1]{\vspace{ #1  cm}}
\newcommand{\compcent}[1]{\vcenter{\hbox{$#1\circ$}}}
\newcommand{\comp}{\mathbin{\mathchoice
{\compcent\scriptstyle}{\compcent\scriptstyle}
{\compcent\scriptscriptstyle}{\compcent\scriptscriptstyle}}}
\newcommand\T{\rule{0pt}{2.6ex}}
\newcommand{\elipse}[2]{\dfrac{x^2}{#1} + \dfrac{y^2}{#2} = 1}
\newcommand{\inner}[2]{\langle \vec{#1}, \vec{#2} \rangle}
\newcommand{\proj}{{\rm proj}}
\newcommand{\eq}{=}
\newcommand{\norm}[1]{\left\Vert{#1}\right\Vert}
\newcommand{\integer}{\mathbb{Z}}
\renewcommand{\sin}{{\rm sen\,}}
\renewcommand{\tan}{{\rm tg\,}}
\renewcommand{\csc}{{\rm cossec\,}}
\renewcommand{\cot}{{\rm cotg\,}}
\renewcommand{\sinh}{{\rm senh\,}}




\title{Quádricas}
\author[\autor]{\autor}
\institute[\instituto]{\instituto}
\date{}

\begin{document}
    \begin{frame}
        \maketitle
    \end{frame}

    \logo{
\includegraphics[scale=.1]{logo-MAT.png}\vspace*{8.5cm}}

    \begin{frame}
      Em $\real^3$ estamos interessados em objetos definidos por uma equação de segundo grau em $x$, $y$ e $z$ da forma\pause
      \begin{equation}\label{equacao_quadrica}
          ax^2 + by^2 + cz^2 + dxy + exz + fyz + gx + hy + iz + j = 0.\pause
      \end{equation}

      O conjunto de pontos em $\real^3$ \pause satisfazendo uma equação da forma \eqref{equacao_quadrica} \pause é chamado de uma \textbf{quádrica} \pause ou \textbf{superfície quádrica}. \pause Nessa equação, \pause pelo menos um dos números \pause $a$, \pause $b$, \pause $c$, \pause $d$, \pause $e$ ou \pause $f$ deve ser não nulo.

    \end{frame}

    \begin{frame}
      \begin{definicao}
          Uma \textbf{esfera} é uma quádrica \pause com equação na forma\pause
          \begin{equation*}
              x^2 + y^2 + z^2 = r^2.
          \end{equation*}
      \end{definicao}
    \end{frame}

    \begin{frame}
      \begin{figure}
          \centering
          \caption{Esfera: $x^2 + y^2 + z^2 = r^2$}
          \includegraphics[scale=0.2]{esfera.png}
      \end{figure}
    \end{frame}

    \begin{frame}
      \begin{definicao}
          Um \textbf{elipsoide} \pause é uma quádrica descrita pela equação\pause
          \begin{equation*}
              \dfrac{x^2}{a^2} + \dfrac{y^2}{b^2} + \dfrac{z^2}{c^2} = 1\pause
          \end{equation*}
          onde os números reais \pause $a$, \pause $b$ e $c$ \pause são positivos \pause e pelo menos dois deles são distintos.
      \end{definicao}
    \end{frame}

    \begin{frame}
      \begin{figure}
          \centering
          \caption{Elipsoide: $\dfrac{x^2}{a^2} + \dfrac{y^2}{b^2} + \dfrac{z^2}{c^2} = 1$}
          \includegraphics[scale=0.2]{elipsoide.png}
      \end{figure}
    \end{frame}

    \begin{frame}
      \begin{definicao}
          Sejam $a$ e $b$ números reais positivos \pause a quádrica descrita de equação\pause
          \[
              z = \dfrac{x^2}{a^2} + \dfrac{y^2}{b^2},\pause
          \]
          é chamada de:\pause
          \begin{enumerate}[label=({\roman*})]
              \item um \textbf{paraboloide elíptico} \pause se $a \ne b$;\pause
              \item um \textbf{paraboloide de rotação} \pause se $a = b$.
          \end{enumerate}
      \end{definicao}
    \end{frame}

    \begin{frame}
      \begin{figure}
          \centering
          \caption{Paraboloide Elíptico: $z = \dfrac{x^2}{a^2} + \dfrac{y^2}{b^2}$}
          \includegraphics[scale=0.2]{paraboloide_eliptico.png}
      \end{figure}
    \end{frame}

    \begin{frame}
        \begin{figure}
            \centering
            \caption{Paraboloide de rotação: $z = \dfrac{x^2}{a^2} + \dfrac{y^2}{a^2}$}
            \includegraphics[scale=0.2]{paraboloide_rotacao.png}
        \end{figure}
    \end{frame}

    \begin{frame}
      \begin{definicao}
          Sejam $a$ e $b$ números reais positivos. \pause A quádrica descrita pela equação\pause
          \[
              z = \dfrac{y^2}{b^2} - \dfrac{x^2}{a^2},\pause
          \]
          é chamada de um um \textbf{paraboloide hiperbólico}.
      \end{definicao}
    \end{frame}

        \begin{frame}
        \begin{figure}
            \centering
            \caption{Paraboloide Hiperbólico: $z = \dfrac{y^2}{b^2} - \dfrac{x^2}{a^2}$}
            \includegraphics[scale=0.2]{paraboloide_hiperbolico.png}
        \end{figure}
    \end{frame}

    \begin{frame}
      \begin{definicao}
          Um \textbf{hiperboloide de uma folha} \pause é uma quádrica descrita pela equação\pause
          \[
              \dfrac{x^2}{a^2} + \dfrac{y^2}{b^2} - \dfrac{z^2}{c^2} = 1\pause
          \]
          onde os números reais $a$, \pause $b$ \pause e $c$ \pause são positivos \pause e pelo menos dois deles são distintos.
      \end{definicao}
    \end{frame}

    \begin{frame}
      \begin{figure}
          \centering
          \caption{Hiperboloide de uma folha: $\dfrac{x^2}{a^2} + \dfrac{y^2}{b^2} - \dfrac{z^2}{c^2} = 1$}
          \includegraphics[scale=0.2]{hiperboloide_de_uma_folha.png}
      \end{figure}
    \end{frame}

    \begin{frame}
      \begin{definicao}
          Um \textbf{hiperboloide de duas folhas} \pause é uma quádrica descrita pela equação\pause
          \[
              \dfrac{z^2}{c^2} - \dfrac{y^2}{b^2} - \dfrac{x^2}{a^2} = 1\pause
          \]
          onde os números reais $a$, \pause $b$ \pause e $c$ \pause são positivos \pause e pelo menos dois \pause deles são distintos.
      \end{definicao}
    \end{frame}

    \begin{frame}
      \begin{figure}
          \centering
          \caption{Hiperboloide duas folhas: $\dfrac{z^2}{c^2} - \dfrac{x^2}{a^2} - \dfrac{y^2}{b^2} = 1$}
          \includegraphics[scale=0.2]{hiperboloide_de_duas_folhas.png}
      \end{figure}
    \end{frame}

    \begin{frame}
      \begin{definicao}
          Um \textbf{cone elíptico} \pause é um conjunto de pontos que satisfaz uma equação da forma\pause
          \begin{equation}\label{cone_eliptico}
              z^2 = \dfrac{x^2}{a^2} + \dfrac{y^2}{b^2}\pause
          \end{equation}
          em que $a$ e $b$ são números reais positivos. \pause Se $a = b$, \pause o cone é chamado de \textbf{cone circular}.
      \end{definicao}
    \end{frame}

    \begin{frame}
      \begin{figure}
          \centering
          \caption{Cone elíptico: $z^2 = \dfrac{x^2}{a^2} + \dfrac{y^2}{b^2}$}
          \includegraphics[scale=0.2]{cone_eliptico.png}
      \end{figure}
    \end{frame}

    \begin{frame}
      Considere uma quádrica de equação
      \begin{equation}\label{equacao_geral_quadrica}
        ax^2 + by^2 + cz^2 + dxy + exz + fyz + gx + hy + iz + j = 0.\pause
      \end{equation}

      Queremos remover os termos $xy$, \pause $xz$ \pause e $yz$ \pause da equação \eqref{equacao_geral_quadrica} \pause para então usar uma translação de eixos, \pause se necessário, \pause e com isso identificar a quádrica definida por tal equação.
    \end{frame}

    \begin{frame}
      Para isso, vamos tentar encontrar uma matriz $Q \pause = [\vec{U_1}\ \pause \vec{U_2}\ \pause \vec{U_3}]$ \pause onde os vetores $\vec{U_1}$, \pause $\vec{U_2}$ \pause e $\vec{U_3}$ \pause são unitários, \pause ortogonais \pause e tais que a mudança\pause
      \begin{equation}\label{substituicao_quadrica}
          \begin{bmatrix}
              x\\
              y\\
              z
          \end{bmatrix} = Q \begin{bmatrix}
              x_1\\
              y_1\\
              z_1
          \end{bmatrix}
      \end{equation}\pause
      transforma a equação \eqref{equacao_geral_quadrica} \pause numa equação da forma\pause
      \begin{equation}\label{equacao_rotacionada_quadrica}
        a'x_1^2 + b'y_1^2 + c'z_1^2 + g'x_1 + h'y_1 + i'z_1 + j = 0.
      \end{equation}
    \end{frame}

    \begin{frame}
      Inicialmente observe que a equação \eqref{equacao_geral_quadrica} \pause pode ser escrita na forma\pause
      \begin{equation}\label{equacao_matricial_quadrica}
        X^tAX + KX + J = [0]\pause
      \end{equation}
      onde
      \[
        A = \begin{bmatrix}
          a & d/2 & e/2\\
          d/2 & b & f/2\\
          e/2 & f/2 & c
        \end{bmatrix}\pause
      , \quad K = \begin{bmatrix}
        g & h & i
      \end{bmatrix}, \pause \quad X = \begin{bmatrix}
        x\\
        y\\
        z
      \end{bmatrix}\pause \quad\mbox{e}\quad J = [j]
      \]\pause
      e $t$ denota a matriz transposta.

    \end{frame}

    \begin{frame}
      Queremos fazer uma substituição da forma \eqref{substituicao_quadrica}, \pause isto é, $X = QX'$ \pause para com isso obter uma equação na forma\pause
      \[
          X_1^tBX_1 + K'X_1 + J = [0]
      \]\pause
      onde\pause
      \[
            B = \begin{bmatrix}
              a' & d'/2 & e'/2\\
              d'/2 & b' & f'/2\\
              e'/2 & f'/2 & c'
              \end{bmatrix} = Q^tAQ,\pause
            K' = \begin{bmatrix}
              g' & h' & i'
              \end{bmatrix} = KQ,\pause
              \]
              \[
            X = \begin{bmatrix}
              x\\y\\z
          \end{bmatrix},\pause
          X_1 = \begin{bmatrix}
              x_1\\y_1\\z_1
          \end{bmatrix}.
        \]
    \end{frame}

    \begin{frame}
      Observe que a matriz $Q$ \pause é formada por vetores unitários \pause e ortogonais, \pause daí $Q^tQ = I_3$ \pause e com isso obtemos\pause
      \[
        B - \lambda I_3 \pause = Q^t A Q - \lambda I_3 \pause = Q^t A Q - Q^t I_3 Q \pause = Q^t (A - \lambda I_3)Q\pause
      \]
      para todo $\lambda \in \real$. \pause Com isso\pause
      \[
        \det(B - \lambda I_3) \pause = \det[Q^t(A - \lambda I_3)Q] = \pause
      \]
      \[
       \det(Q^t)\det(A - \lambda I_3)\det(Q) \pause = \det(A - \lambda I_3).
      \]
    \end{frame}

    \begin{frame}
      Assim escolhendo a matriz $Q$ \pause de tal forma que $d' = e' = f' = 0$, \pause o que é sempre possível, \pause obtemos\pause
      \[
        \det(A - \lambda I_3) \pause = \det(B - \lambda I_3) \pause = \det \begin{bmatrix}
          a' - \lambda & 0 & 0\\
          0 & b' - \lambda & 0\\
          0 & 0 & c' - \lambda
        \end{bmatrix} \pause
        \]
        \[
            = -(\lambda - a')(\lambda - b')(\lambda - c').
        \]
    \end{frame}

    \begin{frame}
      Logo os coeficientes $a'$, $b'$ e $c'$ \pause são raízes da equação de terceiro grau\pause
      \[
        p(\lambda) \pause = \det(A - \lambda I_3) \pause = \det \begin{bmatrix}
          a - \lambda & d/2 & e/2\\
          d/2 & b - \lambda & f/2\\
          e/2 & f/2 & c - \lambda
        \end{bmatrix}.
      \]
    \end{frame}

    \begin{frame}
      Agora para determinar a matriz $Q$ \pause começamos com a equação\pause
      \[
        B = Q^t A Q
      \]\pause
      e multiplicamos à esquerda por $Q$ \pause obtendo\pause
      \[
        Q B = AQ.
      \]
    \end{frame}

    \begin{frame}
      Mas
      \[
          AQ \pause = A \begin{bmatrix}
              \vec{U_1} \pause & \vec{U_2} \pause & \vec{U_3}
          \end{bmatrix} \pause = \begin{bmatrix}
              A\vec{U_1} \pause & A\vec{U_2} \pause & A\vec{U_3}
          \end{bmatrix}
      \]\pause
      e\pause
      \[
          QB \pause = \begin{bmatrix}
              \vec{U_1} \pause & \vec{U_2} \pause & \vec{U_3}
          \end{bmatrix}\pause \begin{bmatrix}
              a' & 0 & 0\\
              0 & b' & 0\\
              0 & 0 & c'
          \end{bmatrix} \pause = \begin{bmatrix}
              a'\vec{U_1} \pause & b'\vec{U_2} \pause & c'\vec{U_3}
          \end{bmatrix}.
      \]
    \end{frame}

    \begin{frame}
      Daí $\vec{U_1}$, \pause $\vec{U_2}$ \pause e $\vec{U_3}$ \pause devem satisfazer as equações\pause
      \[
          A\vec{U_1} \pause = a'\vec{U_1},\pause \quad A\vec{U_2} \pause = b'\vec{U_2} \pause \quad A\vec{U_3} \pause = c'\vec{U_3}.\pause
      \]

      Da primeira equação obtemos\pause
      \begin{align*}
          A\vec{U_1} = a'I_3\vec{U_1}\\
          (A - a'I_3)\vec{U_1} = \vec{0}.
      \end{align*}\pause

      Logo $\vec{U_1}$ \pause é uma solução de norma 1 do sistema linear\pause
      \[
        (A - a'I_3)X = [0].
      \]
    \end{frame}

    \begin{frame}
      Analogamente, $\vec{U_2}$ \pause é uma solução de norma 1 do sistema linear\pause
      \[
        (A - b'I_3)X = [0],\pause
      \]
      que seja ortogonal a $\vec{U_1}$. \pause O terceiro vetor $\vec{U_3}$ \pause é obtido de forma análoga \pause e deve ser ortogonal tanto a $\vec{U_1}$ \pause quanto a $\vec{U_2}$, \pause assim podemos tomar\pause
      \[
          \vec{U_3} = \vec{U_1}\times\vec{U_2}.
      \]
    \end{frame}

    \begin{frame}
      Como $\vec{U_1}$ \pause e $\vec{U_2}$ \pause são de norma 1, \pause então $\vec{U_3}$ \pause também é de norma 1.\pause

      Portanto, a mudança que elemina os termos $xy$, \pause $xz$ \pause e $yz$ \pause é dada por \pause $Q = [\vec{U_1}\  \vec{U_2}\ \vec{U_3}]$. \pause Assim temos o seguinte teorema:
    \end{frame}

    \begin{frame}
      \begin{teorema}
        Considere a equação\pause
        \[
            ax^2 + by^2 + cz^2 + dxy + exz + fyz + gx + hy + iz + j = 0
        \]\pause
        com $a$, $b$, $c$, $d$, $e$, $f$, $g$, $h$, $i$, $j \in \real$, \pause sendo $a$, $b$, $c$, $d$, $e$ e $f$ não necessariamente nulos. \pause Então mediante uma mudança da forma\pause
        \[
          X = Q X_1
        \]\pause
        onde\pause
        \[
          X = \begin{bmatrix}
            x\\y\\z
          \end{bmatrix}, \quad
          Q = \begin{bmatrix}
              \vec{U_1} & \vec{U_2} & \vec{U_3}
        \end{bmatrix}, \quad\mbox{e}\quad
        X_1 = \begin{bmatrix}
          x_1\\y_1\\z_1
        \end{bmatrix}
        \]\pause
        a equação da geral da quádrica \pause acima pode ser transformada em\pause
        \[
          a'x_1^2 + b'y_1^2 + c'z_1^2 + g'x_1 + h'y_1 + i'z_1 + j = 0
        \]
      \end{teorema}
    \end{frame}

    \begin{frame}
      \begin{teorema}
        em que $a'$, \pause $b'$ \pause e $c'$ \pause são raízes de\pause
        \[
          p(\lambda) = \det(A - \lambda I_3) = \det \begin{bmatrix}
            a - \lambda & d/2 & e/2\\
            d/2 & b - \lambda & f/2\\
            e/2 & f/2 & c - \lambda
          \end{bmatrix}.
        \]\pause

        Mais ainda, \pause $\vec{U_1}$ \pause é uma solução de norma 1 \pause do sistema linear\pause
        \[
          \begin{bmatrix}
            a - a' & d/2 & e/2\\
            d/2 & b - a' & f/2\\
            e/2 & f/2 & c - a'
          \end{bmatrix} \begin{bmatrix}
            x\\y\\z
          \end{bmatrix} = \begin{bmatrix}
            0\\0\\0
          \end{bmatrix},
        \]
      \end{teorema}
    \end{frame}

    \begin{frame}
      \begin{teorema}
        $\vec{U_2}$ \pause é uma solução de norma 1 \pause do sistema linear\pause
        \[
          \begin{bmatrix}
            a - b' & d/2 & e/2\\
            d/2 & b - b' & f/2\\
            e/2 & f/2 & c - b'
          \end{bmatrix} \begin{bmatrix}
            x\\y\\z
          \end{bmatrix} = \begin{bmatrix}
            0\\0\\0
          \end{bmatrix},
        \]\pause
        e
        \[
            \vec{U_3} = \vec{U_1} \times \vec{U_2}
        \]
      \end{teorema}
    \end{frame}

    \begin{frame}
      \begin{teorema}
        Seja $\mathcal{S}$ \pause o conjunto dos pontos do espaço \pause que satisfazem a equação\pause
        \[
          ax^2 + by^2 + cz^2 + dxy + exz + fyz + gx + hy + iz + j = 0,\pause
        \]
        com $a$, $b$, $c$, $d$, $e$, $f$, $g$, $h$, $i$, $j \in \real$, \pause sendo $a$, $b$, $c$, $d$, $e$ e $f$ não simultaneamente nulos.
      \end{teorema}
    \end{frame}

    \begin{frame}
      \begin{teorema}
        Sejam $a'$, \pause $b'$ \pause e $c'$ \pause raízes de\pause
        \[
          p(\lambda) = \det \begin{bmatrix}
            a - \lambda & d/2 & e/2 \\
            d/2 & b - \lambda & f/2\\
            e/2 & f/2 & c - \lambda
          \end{bmatrix}.
        \]\pause
        \begin{enumerate}[label=({\roman*})]
          \item Se $a'$, \pause $b'$ \pause e $c'$ \pause tiverem o mesmo sinal, \pause então $\mathcal{S}$ \pause é um elipsoide, \pause um ponto \pause ou o conjunto vazio.\pause
          \item Se $a'$, \pause $b'$ \pause e $c'$ \pause forem não nulos \pause e não tiverem o mesmo sinal, \pause então $\mathcal{S}$ \pause é um hiperboloide de uma folha, \pause de duas folhas \pause ou um cone elíptico.\pause
          \item Se apenas um \pause entre $a'$, \pause $b'$ \pause e $c'$ \pause for nulo, \pause então $\mathcal{S}$ \pause é um paraboloide elíptico, \pause hiperbólico, \pause um cilindro elíptico, \pause hiperbólico, \pause dois planos concorrentes, \pause uma reta \pause ou o conjunto vazio.\pause
          \item Se exatamente dois \pause entre $a'$, \pause $b'$ \pause e $c'$ \pause forem nulos, \pause então $\mathcal{S}$ \pause é um cilindro parabólico, \pause um par de planos paralelos \pause ou um plano.
        \end{enumerate}
      \end{teorema}
    \end{frame}
\end{document}
